\noindent\textbf{Resumen.}\par\smallskip
La selección conjunta excepcional y temporal restringe el dominio de \(19\,446\) candidatos a \(79\) por incidencia y a un único \(K\) por compatibilidad temporal. Se conservan el repertorio, las orientaciones y los filtros que producen esa unicidad.

Se determina la constante holográfica de acoplamiento aritmético--geométrico mediante un registro ordenado \(K\), su evaluación racional y sus representaciones incidenciales, y se demuestra una ley de conservación que permite recuperar información a través de sus transformaciones. Fijadas la base mil y la coordenatización de fase, la lectura \(\kappa=N_K/(1000^{12}-1)\) conserva exactamente las doce componentes del registro. Su órbita cíclica forma una base y proporciona identificación estable de operadores lineales: se utilizan las respuestas a la base completa o una respuesta vectorial cuando existe covarianza cíclica. Así, la estructura que determina el valor participa también en la reconstrucción de las operaciones que actúan sobre ella.

La construcción parte de un núcleo formal denominado Holografía Modular Triádica, en adelante HMT. Comprende \emph{All Possible Paths} (APP), estructura aritmética de caminos sobre un soporte toroidal con grafo de Cayley y evaluaciones aditiva y multiplicativa distinguidas; el TRIT, firma orientada de valores \(\{-1,0,+1\}\); y el \emph{Topological Phase Kernel} (TPK), que compone selección, transporte y memoria. La prolongación nonádica conserva el estado enriquecido del que proceden las coordenadas aritméticas y las incidencias de frontera. El registro \(K\) interviene en la clausura de \(\alpha\) junto con las coordenadas HMT de \(\pi,\varphi,e\); su realización incidencial transporta la misma información ordenada a los espacios que se especifican.

El resultado general de conservación es una identidad bilateral para dos isometrías con dominio y codominio comunes, válida en dimensión arbitraria. El refinamiento entero con acarreo tiene inversa exacta sobre etiquetas admisibles y un levantamiento unitario que se compone con esa identidad. En la especialización nonádica, los pesos \(72/73\) y \(1/73\) determinan la transformación de observables; los invariantes son exactamente los que conmutan con el transporte relativo. Una política explícita de archivo transfiere toda la norma a la memoria infinita y ofrece reconstrucción por un adjunto de norma uno, además de inversas finitas con cotas uniformes en profundidad. Dos observaciones complementarias recuperan con cota óptima el registro centrado, del que se reconstruye el proyector dimensional especificado. La conservación se extiende a cada grado exterior admisible, incluidas las contribuciones mixtas de lectura y memoria.

La extensión reversible relaciona deformación, holonomía y acción; las realizaciones electromagnéticas y gravitatorias conservan sus operadores y normalizaciones explícitos. Para el protocolo simpléctico declarado, con dos entradas gaussianas puras, centradas e idénticas y un número finito de modos, la componente antilineal de la memoria determina el espectro simpléctico reducido, manteniendo todas las ocupaciones de Fock. Un lazo elíptico--parabólico da pureza \(9/11\), espectro \((9/10)10^{-j}\) y entropía adimensional \((10/9)\log10-\log9\). El resultado conjunto caracteriza qué conserva una transformación, qué modifica una observación y qué datos permiten invertirla.

\medskip
% BEGIN RECIPROCIDAD 20260918
La dirección recuperable de \(K\) determina además un flujo positivo de determinante uno, conserva los estratos positroidales y produce redes duales de covolumen fijo con simetría ternaria. Su lectura theta--Mellin conserva polos y residuos, mientras la variación espectral reside en una diferencia entera. El contraste angular generado se realiza como coordenada de Klein; con la orientación conservada, el radio recupera la sección de acción y reexpresa la fórmula electrónica completa. Las pruebas enlazan esa geometría con el archivo bilateral y proporcionan cotas cuantitativas de reconstrucción.

% END RECIPROCIDAD 20260918

% BEGIN R27_FRAMING_SUMMARY_ES
La base emisora de esta recuperación se explicita mediante el alfabeto
de \(42\) bloques decimales, los empalmes regionales y un ejemplo que
separa memoria ordenada e histograma. La comparación de calibres
determina la normalización del observable multiplicativo, y los controles
de clausura separan orientación, condición terminal y posición del registro.
En la frontera conjunta, saturación y neutralidad dual dejan dos
representantes; la orientación fijada determina la sección seleccionada.
Estos resultados identifican los datos que reciben posteriormente
las leyes de conservación y recuperación.
% END R27_FRAMING_SUMMARY_ES

\noindent\textbf{Palabras clave:} Holografía Modular Triádica; Mecánica Dimensional; extrema y media razón holográfica; acoplamiento aritmético--geométrico; refinamiento reversible; memoria operatoria; holonomía; incidencia excepcional; conservación exterior; transporte bilateral; observables invariantes; recuperación estable; acción hamiltoniana; covarianza gaussiana; entropía de entrelazamiento.

\clearpage
\noindent\textbf{Abstract.}\par\smallskip
Joint exceptional and temporal selection restricts the domain of \(19\,446\) candidates to \(79\) through incidence and to a unique \(K\) through temporal compatibility. The repertoire, orientations and filters producing this uniqueness are retained explicitly.

The holographic arithmetic--geometric coupling constant is determined through an ordered register \(K\), its rational evaluation and its incidence representations, and a conservation law is proved that permits information recovery through these transformations. With base one thousand and the phase chart fixed, the reading \(\kappa=N_K/(1000^{12}-1)\) preserves all twelve register components exactly. Its cyclic orbit forms a basis and provides stable identification of linear operators: the data are the responses to the complete basis, or one vector response under cyclic covariance. The structure that determines the value thus also participates in reconstructing the operations acting on it.

The construction starts from a formal kernel called Triadic Modular Holography, hereafter HMT. It comprises \emph{All Possible Paths} (APP), an arithmetic path structure on a toroidal support with a Cayley graph and distinct additive and multiplicative evaluations; TRIT, an oriented signature taking values in \(\{-1,0,+1\}\); and the \emph{Topological Phase Kernel} (TPK), which composes selection, transport and memory. Nonadic prolongation preserves the enriched state from which arithmetic coordinates and boundary incidences arise. The register \(K\) enters the closure of \(\alpha\) together with the HMT coordinates of \(\pi,\varphi,e\); its incidence realization transports the same ordered information to the specified spaces.

The general conservation result is a bilateral identity for two isometries with a common domain and codomain, valid in arbitrary dimension. Integer refinement with carry has an exact inverse on admissible labels and a unitary lift that composes with this identity. In the nonadic specialization, the weights \(72/73\) and \(1/73\) determine the transformation of observables; the invariants are precisely those commuting with the relative transport. An explicit archive policy transfers the entire norm to infinite memory and yields reconstruction by an adjoint of norm one, together with finite inverses whose bounds are uniform in depth. Two complementary observations recover the centred register with an optimal bound, from which the specified dimensional projector is reconstructed. Conservation extends to every admissible exterior degree, including mixed readout--memory contributions.

The reversible extension relates deformation, holonomy, and action; the electromagnetic and gravitational realizations retain their explicit operators and normalizations. For the declared symplectic protocol with two identical centred pure Gaussian inputs and finitely many modes, the antilinear component of memory determines the reduced symplectic spectrum, retaining all Fock occupations. An elliptic--parabolic loop gives purity \(9/11\), spectrum \((9/10)10^{-j}\), and dimensionless entropy \((10/9)\log10-\log9\). Taken together, the results characterize what a transformation preserves, what an observation changes, and which data permit its inversion.

\medskip
% BEGIN RECIPROCIDAD 20260918
The recoverable direction of \(K\) also determines a positive determinant-one flow, preserves positroid strata, and produces dual lattices of fixed covolume with ternary symmetry. Their theta--Mellin readout preserves poles and residues, while spectral variation lies in an entire difference. The generated angular contrast is realized as a Klein coordinate; with orientation retained, the radius recovers the action section and re-expresses the complete electronic formula. The proofs connect this geometry with the bilateral archive and provide quantitative reconstruction bounds.

% END RECIPROCIDAD 20260918

% BEGIN R27_FRAMING_SUMMARY_EN
The emitting basis of this recovery is made explicit through the alphabet
of \(42\) decimal blocks, regional junctions, and an example separating
ordered memory from its histogram. Gauge comparison determines the
normalization of the multiplicative observable, while closure controls
separate orientation, terminal condition, and register position. At the
joint boundary, saturation and dual neutrality leave two representatives;
the fixed orientation determines the selected section. These results
identify the data subsequently received by the conservation and
recovery laws.
% END R27_FRAMING_SUMMARY_EN

\noindent\textbf{Keywords:} Triadic Modular Holography; Dimensional Mechanics; holographic extreme and mean ratio; arithmetic--geometric coupling; reversible refinement; operator memory; holonomy; exceptional incidence; exterior conservation; bilateral transport; invariant observables; stable recovery; Hamiltonian action; Gaussian covariance; entanglement entropy.
